% SPDX-License-Identifier: CC-BY-SA-4.0
% Author: Matthieu Perrin

\documentclass[livret]{td}

\usepackage{src/courses/pcmt/config}

\typeTD



\hypersetup{
  pdftitle  = {Programmation Concurrente en Multi-Threads},
  pdfauthor = {Matthieu Perrin},
  pdfsubject  = {TD de M1 du cours Programmation Concurrente en Multi-Threads},
  pdfkeywords = {concurrence, Java, modèle de mémoire, moniteur, parallélisme, sûreté, thread, verrou, vivacité}
}

\begin{document}

\maketitlepage

\session{Du parallélisme à la concurrence}
% SPDX-License-Identifier: CC-BY-SA-4.0
% Session: Form parallelism to concurrency
% Exercise: Amdahl's law

\begingroup

\begin{exercice}[La loi d'Amdahl]
  \label{exo:introduction/amdahl}

  Une proportion $p=70\,\%$ d'un programme séquentiel peut être parallélisée. On dispose d'une machine avec $n=10$~c\oe urs.

  \begin{question} 
  \item Quelle est l'accélération (speedup) attendue ?
  \end{question}
  \begin{correction}
    $$ S(10)=\dfrac{1}{(1-p)+\dfrac{p}{10}} = \dfrac{1}{0,3+\dfrac{0,7}{10}} \approx 2,70. $$
  \end{correction}

  \begin{question} 
  \item On trouve dans le commerce une machine avec $n=20$~c\oe urs. Quelle est l'accélération attendue ?
  \end{question}
  \begin{correction}
    $$ S(20)=\dfrac{1}{(1-p)+\dfrac{p}{20}} = \dfrac{1}{0,3+\dfrac{0,7}{20}} \approx 2,99. $$
  \end{correction}

  \begin{question} 
  \item Quelle est l'accélération maximale théorique lorsque $n \to \infty$ ?
  \end{question}
  \begin{correction}
    $$ \lim_{n\to\infty} S(n) = \dfrac{1}{1-p} = \dfrac{1}{0,3} \approx 3,33. $$
  \end{correction}

  Un ingénieur pense qu'il existe un moyen d'optimiser la partie séquentielle pour la rendre deux fois plus rapide.

  \begin{question} 
  \item Quelle proportion du programme est parallélisable dans la version otimisée ? 
  \end{question}
  \begin{correction}
    On suppose maintenant que la partie séquentielle est deux fois plus rapide,
    mais que la partie non parallélisable est inchangée.
    On obtient donc : 
    $$T_1^{\mathrm{opt}} = \dfrac{T_1 \times (1-p)}{2} + T_1 \times p. $$

    La proportion parallélisable est donc :
    $$ p_{\mathrm{opt}} = \dfrac{T_1 \times p}{T_1^{\mathrm{opt}}} = \dfrac{2p}{1+p} \approx 82\%. $$
  \end{correction}
  
  \begin{question} 
  \item Quelle accélération obtient-on alors, dans la parallélisation avec $n = 10$~c\oe urs ?
  \end{question}
  \begin{correction}
    $$ S_{\mathrm{opt}}(10)=\dfrac{1}{(1-p_{\mathrm{opt}})+\dfrac{p_{\mathrm{opt}}}{10}} \approx 3,86. $$
  \end{correction}

  \begin{question} 
  \item À coût égal, vaut-il mieux choisir d'acheter la machine avec $n=20$~c\oe urs ou payer l'ingénieur pour optimiser le code ? 
  \end{question}
  \begin{correction}
    Il faut comparer l'accélération par rapport à $T_{1}$,
    de $T_{20}$, le temps d'exécution du programme initial sur la machine à $20$~c\oe urs,
    et de $T_{10}^{\mathrm{opt}}$, le temps d'exécution du programme optimisé sur la machine à $10$~c\oe urs.
    \begin{itemize}
    \item $\dfrac{T_{20}}{T_1} = S(20) \approx 2,99.$
    \item $\dfrac{T^{\mathrm{opt}}_{10}}{T_1} = \dfrac{1-p}{2} + \dfrac{p}{10} \approx 4,55.$
    \end{itemize}    

    L'optimisation de la partie séquentielle est nettement plus rentable.
  \end{correction}
  
\end{exercice}

\endgroup
\endinput

% SPDX-License-Identifier: CC-BY-SA-4.0
% Session: Form parallelism to concurrency
% Exercise: Parallelization of matrix multiplication

\begingroup

\begin{exercice}[Parallélisation]
  \label{exo:introduction/matrix}

  Le but de cet exercice est d'écrire une fonction \lstinline{double[][] product(double[][] M1, double[][] M2)}
  qui calcule le produit de deux matrices carrées $M_1$ et $M_2$ de taille $n\times n$.
  On rappelle la définition de la multiplication des matrices. Si $M_3 = M_1 \times M_2$, alors pour tout
  $i,j \in \{0,\dots,n-1\}$,
  $$ M_3[i][j] = \sum_{k=0}^{n-1} M_1[i][k] \times M_2[k][j].$$

  \begin{question}
  \item Implémentez la version séquentielle de \lstinline{product}.
  \end{question}
  \begin{correction}
    \begin{lstlisting}[gobble=6]
      public interface Product {
	double[][] product(double[][] M1, double[][] M2) throws InterruptedException;
      }
      public final class Sequential implements Product {
	double[][] product(double[][] M1, double[][] M2) {
	  int n = M1.length;
	  double[][] M3 = new double[n][n];
	  for (int i = 0; i < n; i++) {
	    for (int j = 0; j < n; j++) {
	      M3[ii][jj] = 0;
	      for (int k = 0; k < n; k++) {
		M3[ii][jj] += M1[ii][k] * M2[k][jj];
	      }
	    }
	  }
          return M3;
        }
      }
    \end{lstlisting}
  \end{correction}

  \begin{question}
  \item Proposez une parallélisation où chaque cellule $M_3[i][j]$ est calculée par un thread distinct (jusqu'à $n^2$ threads). 
  \end{question}
  \begin{correction}
    \begin{lstlisting}[gobble=6]
      package td1_parallelism.matrix;

      public final class Naive implements Product {

	public double[][] product(double[][] M1, double[][] M2) throws InterruptedException {

	  int n = M1.length;
	  double[][] M3 = new double[n][n];
	  Thread[] threads = new Thread[n*n];

	  for (int i = 0; i < n; i++) {
	    for (int j = 0; j < n; j++) {
	      final int ii = i, jj = j;
	      threads[i*n+j] = new Thread(() -> {
		M3[ii][jj] = 0;
		for (int k = 0; k < n; k++) {
		  M3[ii][jj] += M1[ii][k] * M2[k][jj];
		}
	      });
	    }
	  }

	  for(var t : threads) t.start();
	  for(var t : threads) t.join();

	  return M3;
	}

      }
    \end{lstlisting}
  \end{correction}

  \begin{question}
  \item Proposez une version qui crée au plus $T$ threads ($T$ est une variable globale), chacun calculant un lot de cellules.
    Précisez votre stratégie de répartition.
  \end{question}
  \begin{correction}
    \begin{lstlisting}[gobble=6]
      public static double[][] product(int n, double[][] M1, double[][] M2) {
        double[][] M = new double[n][n];
        Thread threads[][] = new Thread[n][n];
        for(int i = 0; i < n; i++) {
	  int thisI = i;
	  threads[i][0] = new Thread(() -> {
	    for(int j = 0; j < n; j++) {
	      int thisJ = j;
	      threads[thisI][j] = new Thread(() -> {
	        M[thisI][thisJ] = 0;
	        for(int k = 0; k<n; k++) {
		  M[thisI][thisJ]+=M1[thisI][k] * M2[k][thisJ];
	        }
	      });
	      threads[thisI][j].start();
	    }
	    for(int j = 1; j < n; j++) {
	      try {
                threads[thisI][j].join();
              } catch (InterruptedException e) { }
	    }
	    
	  });
	  threads[i][0].start();
        }
        for(int i = 1; i < n; i++) {
	  try {
            threads[i][0].join();
          } catch (InterruptedException e) { }
        }
        return M;
      }
    \end{lstlisting}
  \end{correction}

  \begin{question}
  \item Analysez la complexité en temps des versions et justifiez l'intérêt
    de limiter le nombre de threads à une borne indépendante de $n$.
  \end{question}
  \begin{correction}
    Si on parallélise la boucle intérieure, on passe d'une complexité $O(n^3)$ à $O(n^2)$, soit une accélération en $O(n)$.
    On peut faire mieux en parallélisant la boucle extérieure également. On passe alors de $O(n^3)$ à $O(n)$, soit une accélération en $O(n^2)$.
    Cela ne viole pas la loi d'Amdahl car on change le problème en changeant le nombre de threads :
    on augmente le nombre de threads pour s'adapter à la quantité de données à traîter. C'est l'idée de la loi de Gustafson :
    Si le nombre de threads est proportionnel à la taille des données, on peut espérer une accélération en $O(n)$.
    
    Remarque : le programme ci-dessus utilise $n^2 + n + 1$ threads, mais on peut utiliser les mêmes threads pour paralléliser et pour faire certains calculs.
  \end{correction}
  
\end{exercice}

\endgroup
\endinput

% SPDX-License-Identifier: CC-BY-SA-4.0
% Session: Form parallelism to concurrency
% Exercise: Asynchrony

\begingroup

\begin{exercice}[Asynchronie]
  \label{exo:introduction/asynchronous}

  On se donne le programme multi-threads suivant : 

  \begin{lstlisting}
    Thread t = new Thread(() -> {
      System.out.print("a");
      new Thread(() -> { System.out.print("b"); }).start();
      System.out.print("c");
    });

    System.out.print("d");
    t.start();
    System.out.print("e");
    t.join();
    System.out.print("f");
  \end{lstlisting}

  \begin{question}
  \item Décrivez la relation ``happens-before'' du programme.
  \item Décrivez toutes les sorties possibles pour le programme, en sachant qu'il est séquentiellement cohérent.
  \item Décrivez la machine à états du programme.
  \end{question}

  \begin{correction}
    Les sorties possibles sont toutes les linéarisations du graphe de causalité, soit les 11 sorties suivantes :
    
    \begin{minipage}{.4\textwidth}
      \begin{tikzpicture}
        \draw (0,0) node {f};
        \draw (1,1) node {c};
        \draw (3,1) node {b};
        \draw (0,2) node {e};
        \draw (2,2) node {a};
        \draw (1,3) node {d};
        
        \draw[-latex] (0.8,2.8) -- (0.2,2.2);
        \draw[-latex] (1.2,2.8) -- (1.8,2.2);
        \draw[-latex] (0,1.8) -- (0,0.2);
        \draw[-latex] (1.8,1.8) -- (1.2,1.2);
        \draw[-latex] (2.2,1.8) -- (2.8,1.2);
        \draw[-latex] (0.8,0.8) -- (0.2,0.2);
      \end{tikzpicture}
    \end{minipage}
    \begin{minipage}{.4\textwidth}
      \begin{itemize}
      \item deacfb
      \item deacbf
      \item deabcf
      \item daecfb
      \item daecbf
      \item daebcf
      \item dabecf
      \item dacefb
      \item dacebf
      \item dacbef
      \item dabcef
      \end{itemize}
    \end{minipage}
  \end{correction}
  
\end{exercice}

\endgroup
\endinput

\sessionappendix
% SPDX-License-Identifier: CC-BY-SA-4.0
% Session: Form parallelism to concurrency
% Exercise: The mad scientist

\begingroup

\begin{exercice}[Le savant fou%
    \footnote{Inspiration : M. Herlihy \& N. Shavit. \emph{The Art of Multiprocessor Programming}. Morgan Kaufmann (2008)}]
  \label{exo:introduction/mad_scientist}

  En visite pédagogique dans un laboratoire, vous et le reste de votre groupe êtes 
  capturés par un savant fou. Il fait l'annonce suivante :

  \og
  J'ai mis en place une pièce qui contient une lumière et un interrupteur
  qui la commande. Actuellement, la lumière est éteinte.
  Vous pouvez discuter ce soir d'une stratégie commune, mais à partir de demain
  vous serez isolés et ne pourrez plus communiquer.
  Chaque jour, je tirerai au hasard de manière uniforme l'un d'entre vous, qui passera la journée dans la pièce.
  À tout moment durant sa visite, il pourra déclarer~: \emph{Nous avons tous visité la pièce}.
  S'il a raison, vous êtes libres. Sinon, je vous reboote les cerveaux !
  \fg

  Hypothèses implicites :
  \begin{itemize}
  \item La lumière est éteinte au départ.
  \item L'interrupteur est binaire (allumé/éteint).
  \item La lumière est invisible depuis l'extérieur de la pièce,
    et il n'y a aucun autre moyen de communication entre les étudiants.
  \end{itemize}
  
  \begin{question}
  \item Énoncez précisément ce que doit garantir une stratégie réussie.
  \end{question}
  \begin{correction}
    Une stratégie correcte doit vérifier :
    \begin{description}
    \item[Sûreté.] Si un étudiant déclare \og Nous avons tous visité la pièce \fg,
      alors tous les étudiants l'ont effectivement visitée.
    \item[Vivacité.] Avec probabilité 1, un étudiant finira par pouvoir faire cette déclaration.
    \end{description}
  \end{correction}
  
  \begin{question}
  \item Proposez une stratégie et justifiez brièvement en quoi elle satisfait l'exigence précédente.
  \end{question}
  \begin{correction}
    Une stratégie possible consiste à désigner un étudiant comme \emph{délégué} :
    \begin{itemize}
    \item Chaque étudiant non délégué allume la lumière la première fois qu'il entre
      dans la pièce et la trouve éteinte, puis ne la rallume plus jamais.
    \item Chaque fois que le délégué entre dans la pièce et voit la lumière allumée,
      il l'éteint et incrémente un compteur local.
    \item Lorsque son compteur atteint $n-1$, il déclare que tous les étudiants ont visité la pièce.
    \end{itemize}
    Cette stratégie garantit la sûreté (le délégué ne peut compter qu'un allumage par étudiant)
    et la vivacité presque sûrement (le tirage aléatoire finit par sélectionner chaque étudiant
    au moins une fois alors que la lumière est éteinte).
  \end{correction}
  
  \begin{question}
  \item Reprenez la question précédente, sans l'hypothèse sur l'état initial de la lumière.
  \end{question}
  \begin{correction}
    Si l'état initial de la lumière est inconnu,
    le protocole peut être adapté en demandant au délégué d'ignorer le premier allumage observé
    (pour éliminer l'incertitude sur l'état de départ).
  \end{correction}
  
  \begin{question}
  \item Reprenez les deux premières questions, sans l'hypothèse que le tirage est uniforme. À la place, on suppose
    que si personne ne déclare que tout le monde a visité la pièce, chaque étudiant est sélectionné une infinité de fois
    (pas nécessairement à intervalles bornés). 
  \end{question}
  \begin{correction}
    Si le tirage n'est pas uniforme mais seulement équitable
    (chaque étudiant est choisi une infinité de fois, sans borne sur l'intervalle entre deux visites),
    la stratégie reste correcte : la sûreté n'est pas affectée,
    et la vivacité est renforcée et devient déterministe : un étudiant finira par pouvoir faire cette déclaration.
  \end{correction}

\end{exercice}

\endgroup
\endinput


\session{Synchronisation par verrous}
\input{blocking/philosophers}
\input{blocking/bank}
\input{blocking/pipeline}
\sessionappendix
\input{blocking/tree}

\session{Moniteurs de synchronisation}
\input{monitors/stack}
\input{monitors/unisex}
\input{monitors/threadpool}
\sessionappendix
\input{monitors/barrier}
\input{monitors/savages}
\input{monitors/go}

\session{Le modèle mémoire de Java}
\input{volatile/volatile}
\input{volatile/singleton}
\newpage
% SPDX-License-Identifier: CC-BY-SA-4.0
% Session: Java Memory Model
% Exercise: Identify data races

\begingroup

\begin{exercice}[Identification de \emph{data races}]
  \label{exo:memory/dataraces}

  On examine le programme Java suivant. Les trois variables partagées sont initialisées à \lstinline{x = 0} et \lstinline{flag = true}.
  On suppose qu'aucun autre thread n'intervient et que les réveils intempestifs possibles de \lstinline{wait()} sont correctement
  gérés par la boucle \lstinline{while}. On raisonne sous l'hypothèse de \emph{cohérence séquentielle} pour les questions 1–3.

  \begin{lstlisting}[numbers=left]
    final class Program {
      static final Object o = new Object();
      static int x = 0;
      static boolean flag = true;

      public static void main(String[] args) throws Exception {
        Thread t = new Thread(() -> {
          x = 1;                 
          synchronized (o) {     
            flag = false;        
            o.notify();          
          }                      
          x = 2;                 
        });

        synchronized (o) {       
          t.start();             
          while (flag) {         
            o.wait();            
          }
          int a = x;             
        }                        
        int b = x;               
      }
    }
  \end{lstlisting}

  \begin{enumerate}
  \item Donnez la machine à états du programme.
  \item Déduisez-en toutes les exécutions séquentiellement cohérentes. Quelles sont les valeurs possibles de $a$ et $b$ ?
  \item Identifiez toutes les \emph{data races} présentes dans le programme.
  \item Donnez deux corrections \emph{minimales} et \emph{distinctes} qui éliminent toutes les \emph{data races} :
    \begin{itemize}
    \item en ne modifiant que des qualificatifs (ex. \lstinline{volatile}) dans les déclarations ;
    \item en ne modifiant que la synchronisation (ex. blocs \lstinline{synchronized}).
    \end{itemize}
  \item Pour chacune des deux corrections, quelles sont alors les valeurs possibles de $a$ et $b$ ?
  \end{enumerate}

\end{exercice}

\endgroup
\endinput


\session{Algorithmes d'exclusion mutuelle}
\input{locks/filter}
\input{locks/reentrance}
\input{locks/getAndIncrement}
\sessionappendix
\input{locks/readwrite}

\end{document}

\endinput
